\documentclass[10pt,a4paper]{amsart}
\usepackage[margin=19mm]{geometry}
\usepackage{amsmath,amssymb,mathtools}
\usepackage[scr=rsfs,cal=euler]{mathalfa}
\usepackage{xfrac}
\usepackage[unicode,hidelinks,bookmarks=false]{hyperref}
\newcommand{\ie}{i.e.\ }
\newcommand{\cf}{cf.\ }
\newcommand{\resp}{respectively\ }
\newcommand{\Subsection}{\subsection}
\providecommand{\nobreakdash}{\nobreak\hbox{-}}
\setlength{\emergencystretch}{2em}
\multlinegap0pt
\usepackage{bm}
\usepackage{bbm}
\usepackage{dsfont}
\usepackage{xspace}
\usepackage{cancel}
\usepackage{mathtools}
\usepackage{amsthm}
\makeatletter
\@ifundefined{theorem}{\newtheorem{theorem}{Theorem}}{}
\@ifundefined{corollary}{\newtheorem{corollary}[theorem]{Corollary}}{}
\@ifundefined{definition}{\newtheorem{definition}[theorem]{Definition}}{}
\@ifundefined{lemma}{\newtheorem{lemma}[theorem]{Lemma}}{}
\@ifundefined{proposition}{\newtheorem{proposition}[theorem]{Proposition}}{}
\@ifundefined{remark}{\newtheorem{remark}[theorem]{Remark}}{}
\makeatother
\newcommand{\RR}{{\mathbb R}}
\renewcommand{\tilde}[1]{\widetilde{#1}}
\title[On the dimensions of correlated equilibrium polytopes of gen]{On the dimensions of correlated equilibrium polytopes of generic games}
\author{Jan Draisma, Linda Hoyer, Irem Portakal}
\date{Source: arXiv:2608.04924v1 (CC BY 4.0)}
\begin{document}
\maketitle

\noindent\emph{Source and licence.} On the dimensions of correlated equilibrium polytopes of generic games. Version arXiv:2608.04924v1 (CC BY 4.0), distributed under the Creative Commons Attribution 4.0 International licence, as recorded in the arXiv licence metadata for this exact version and shown on the abstract page https://arxiv.org/abs/2608.04924v1. Full rights notice: licenses/arxiv-2608.04924-CC-BY-4.0.txt.\par\smallskip

\noindent\textit{Editorial scope.} Counted unit: the Theorem whose source label is \texttt{TheoremFullOrSingleton} in the article (source0.tex lines 568--571, source label \texttt{TheoremFullOrSingleton}), delivered together with the article's own complete original proof of it (source0.tex lines 572--684, one \texttt{proof} environment of 113 lines). No mathematics is written, repaired or re-derived: statement and proof are the article's own text line for line. Required context retained uncounted: re:IncConsEq (lines 183--189); PropositionBalancing (lines 342--347); de:Partition (lines 447--461); eq:rp (lines 457--458); RemarkR (lines 462--493); LemmaSubgameAdd (lines 495--504); CorollaryNonZero (lines 521--527); prop:NotAllNash (lines 550--552). Every definition, display and label of those lines is retained unchanged. Omitted from this package and not redistributed: 1--182, 190--341, 348--446, 459--461, 494--494, 505--520, 528--549, 553--567, 685--885. Nothing inside a retained excerpt is omitted or altered: every retained body is delivered exactly as the article's lines are written, so no omission receipt of any kind accompanies this package. Citations: the retained excerpts carry no citation command at all, so no bibliography entry of the article is transcribed and no reference is redistributed. Reference adaptation: none was needed and none was made; every reference inside the delivered unit resolves inside it, so no unresolved reference or citation is produced. Printed slips kept literal: no printed word, symbol, hypothesis or number of the counted statement or of its proof was altered. Layout and class: the article's own class is \texttt{amsart} with packages including geometry, xcolor, appendix, fontenc, amsthm, amsmath, hyperref, comment, caption, graphicx, subfigure, latexsym, amssymb, float; the wrapper is the lane's pinned \texttt{amsart} editorial wrapper at 10pt on A4, which restates the theorem-like environments the retained text needs and adds the article's own macro definitions that the retained text needs (\texttt{\textbackslash RR}, \texttt{\textbackslash tilde}), with extra wrapper packages (bm, bbm, dsfont, xspace, cancel, mathtools). The delivered numbering is the unit's own. Assets: no retained span contains a figure environment, a graphics-inclusion command, a PDF-page inclusion, a table or a TikZ picture; no image file is copied into this package, the package declares no asset dependency and no image file is redistributed with it. Licence: version arXiv:2608.04924v1 of this article is distributed under the Creative Commons Attribution 4.0 International licence, as recorded in the arXiv licence metadata for this exact version and shown on the abstract page https://arxiv.org/abs/2608.04924v1. A full rights notice is delivered at licenses/arxiv-2608.04924-CC-BY-4.0.txt. Characters: every retained excerpt is pure ASCII, so the delivered engine, XeLaTeX with its default Latin Modern text font, reports no missing character.
\begin{remark} \label{re:IncConsEq}
For future use, we make the following observation: if $p \in P_G$, $i \in
[n]$, and $k,\ell \in [d_i]$ are such that the slices $p_{-i,k}$ and $p_{-i,\ell}$
are both non-zero and scalar multiples of each other, then in fact the
incentive constraints $H^{X^{(i)}}_{k,\ell}(p_{-i,k}) \geq 0$ and
$H^{X^{(i)}}_{\ell,k}(p_{-i,\ell}) \geq 0$ both hold with equality.
\end{remark}

\begin{proposition}
\label{PropositionBalancing}
Let $G=(n,S,X)$ and let $p \in P_G$. Then there
exist $Y$ arbitrarily close to $X$ such that the game $G':=(n,S,Y)$
satisfies $p \in P_{G'}$ and such that $p$ is balanced for $G'$.
\end{proposition}

\begin{definition} \label{de:Partition}
Let $p \in \Delta_{D-1}$. For each $i \in [n]$, we define an
equivalence relation on $[d_i]$ where $k \sim \ell$ if and only if
$\pi_{-i,k}(p)=\pi_{-i,\ell}(p)$. We let \[M^{(i)}_1 \dot{\cup}
\dots \dot{\cup} M^{(i)}_{m_i} = [d_i]\] be the associated set
partition and set $d^{(i)}_a:=|M^{(i)}_a|$ for $a \in [m_i]$.



\noindent We define 
\begin{equation} \label{eq:rp} r(p):=\sum_{i=1}^n |\{ \pi_{-i,k}(p)
\mid p_{-i,k} \neq 0, k \in [d_i]   \}|; \end{equation}
note that the $i$-th term here is precisely the cardinality of the set
$W$ in the proof of Proposition~\ref{PropositionBalancing}.
\end{definition}

\begin{remark}
\label{RemarkR}
Let $p \in \Delta_{D-1}$.
\begin{enumerate}
    \item It holds that \[r(p) = \left( \sum_{i=1}^n m_i \right) - |\{i \in [n] \mid p_{-i,k}=0 \ \text{for some} \ k \in [d_i]   \}  |.  \]
    \item It holds that $n \leq r(p) \leq \sum_{i=1}^n d_i$, with
    $r(p)=n$ if and only if $p=(x^{(1)} \otimes \dots \otimes
    x^{(n)})$, where $(x^{(1)}, \dots, x^{(n)})$ is a Nash equilibrium
    of $G$. Indeed, in the latter case, all non-zero slices of $p$ in
    the $i$-th direction are positive scalar multiples of $x_1 \otimes
    \cdots \otimes x_{i-1} \otimes x_{i+1} \otimes \cdots \otimes x_n$, so each $i \in [n]$
    contributes $1$ to the sum in \eqref{eq:rp}. Conversely, if each $i$
    contributes $1$, then it follows that in each direction, the
    non-zero slices of $p$ are positive scalar multiples of each other,
    so that $p$ has rank one as claimed. 
    \item Let $q \in \Delta_{D-1}$ and let $p':=(1- \varepsilon)p +
    \varepsilon q \in \Delta_{D-1}$ for some $0< \varepsilon <1$.
    Then for $\varepsilon$ small enough, if the $k$-th and $\ell$-th
    slice of $p$ in the $i$-th direction are linearly independent, then 
    the same holds for $p'$, hence $r(p') \geq r(p)$. 
    \item Let $a_1 \in [m_1], \dots, a_n \in [m_n]$ be indices and
    consider \[ \tilde{p}:=(p_{j_1, \dots, j_n})_{j_i \in
    M^{(i)}_{a_i}, i \in [n]}.\] Then $\tilde{p}=0$ if 
    for at least one $i \in [n]$, we have $\pi_{-i,k}(p)=0$ for some,
    and hence all, $k \in M_{a_i}^{(i)}$. If that is not the case,
    then in each direction all slices of $\tilde{p}$ 
    are positive scalar multiples of each other. This implies
    that $\tilde{p}=y^{(1)} \otimes \cdots \otimes y^{(n)}$ for
    certain vectors $y^{(i)} \in \RR^{M^{(i)}_{a_i}}$ with 
    positive entries. 
\end{enumerate}
\end{remark}

\begin{lemma}
\label{LemmaSubgameAdd}
Let $p \in P_G$ be balanced. Let $a_1 \in [m_1], \dots, a_n \in [m_n]$ such that $p_{-i,k} \neq 0$ for any $i \in [n]$, $k \in M^{(i)}_{a_i}$ and
let $\tilde{G}=(n,\tilde{S},\tilde{X})$ be the subgame of $G$
corresponding to the subsets $M_{a_i}^{(i)} \subseteq S^{(i)}$. Let
$\tilde{q} \in P_{\tilde{G}}$ be a correlated equilibrium, and extend
$\tilde{q}$ to a tensor $q \in \Delta_{D-1}$ by assigning $0$ to any
tuple of pure strategies not in $\tilde{S}$. Then there exists $0<\varepsilon<1$
such that $(1-\varepsilon)p + \varepsilon q \in P_G$.
\end{lemma}

\begin{corollary}
\label{CorollaryNonZero}
Assume there exists a $p \in P_G$ such that $p_{-i,k} \neq 0$ for any
$i \in [n], k \in [d_i]$. Then there exist $X'$ arbitrarily close to $X$
such that the game $G':=(n,S,X')$ has a $p' \in P_{G'}$ all of whose
entries are positive and that satisfies $r(p') \geq r(p)$.
\end{corollary}

\begin{proposition} \label{prop:NotAllNash}
    Let $G=(n,S,X)$ be a generic game. If $P_G$ is a not a singleton, then it contains points $p \in \Delta_{D-1}$ that are {\em not} of the form $x_1 \otimes \cdots \otimes x_n$ for $(x_1,\ldots,x_n)$ a Nash equilibrium of $G$.
\end{proposition}

\begin{theorem}
\label{TheoremFullOrSingleton}
Let $G=(n,S,X)$ be a generic game. Assume there exists a $p \in P_G$ such that $p_{-i,k} \neq 0$ for any $i \in [n]$, $k \in [d_i]$. Then the correlated equilibrium polytope $P_G$ is either full-dimensional or a singleton.     
\end{theorem}

\begin{proof}
Assume $P_G$ is not a singleton. By Proposition~\ref{prop:NotAllNash}, we may choose $p \in P_G$ with $p_{-i,k} \neq 0$ for all $i \in [n]$, $k \in [d_i]$ that moreover is not a Nash equilibrium, so that $n<r(p)$. 
We argue by induction on the number \[
\left(\sum_{i=1}^n d_i\right)-r(p).\] 
If $r(p)=\sum_{i=1}^n d_i$, then all
incentive constraints for $p$ are satisfied with a strict inequality
and thus $P_G$ is full-dimensional. So assume from now on that that $r(p)<\sum_{i=1}^n d_i$. 

By Proposition \ref{PropositionBalancing} and Corollary
\ref{CorollaryNonZero}, the point $p$ can be assumed balanced and to
have only positive entries. 

We will now construct
a rank-1-tensor $q \in \Delta_{D-1}$, together with an $0<\varepsilon<1$
and an $X'$ arbitrarily close to $X$ such that for the game
$G':=(n,S,X')$ we have $p':=(1-\varepsilon)p
+ \varepsilon q \in P_{G'}$ and $r(p')>r(p).$ By
the induction hypothesis, the theorem then follows.

To find $q$, we begin by relabeling the players, their pure strategies, and the
parts in the partitions of the $[d_i]$ conveniently, as follows:
\begin{enumerate}
    \item Assume that $M^{(i)}_1=[n_1^{(i)}]$ for all $i \in [n]$. 
    \item Since $r(p)< \sum_{i=1}^n d_i$, there is at least one $i$
    for which the partition of $[d_i]$ from
    Definition~\ref{de:Partition} has a part that is not a
    singleton. So we may assume that $n_1^{(1)} \geq 2$.
    \item The condition that $r(p)>n$ implies that, as a tensor, $p$
    has rank $>1$. Then it has some flattening to a matrix that also
    has rank $>1$. This implies that there are at least two indices
    $i$ with $m_{i}>1$, say $i_1$ and $i_2$. 
\end{enumerate}

Let $\tilde{G}=(n,\tilde{S},\tilde{X})$ be the subgame with
$\tilde{S}:=\prod_{i=1}^n M^{(i)}_1$. Let $(y^{(1)},\dots, y^{(n)})$
be a Nash equilibrium of $\tilde{G}$ and let $\tilde{q}:=y^{(1)}
\otimes \dots \otimes y^{(n)} \in P_{\tilde{G}}$, which we extend to a
tensor $q \in \Delta_{D-1}$ by assigning $0$ to any tuple of pure
strategies not in $\tilde{S}$. 

For each $i \in [n]$, the slices $p_{-i,k}$ with $k \in [d^{(i)}_1]$
are all non-zero scalar multiples of each other. Hence 
there is a unique $x^{(i)} \in \Delta_{d^{(i)}_1 -1}$ such that
\[p_{-i,k}= x^{(i)}_k  \sum_{\ell =1}^{d^{(i)}_1} p_{-i,\ell}\]
for those $k$. We define the tensor $\tilde{t} := x^{(1)} \otimes \dots \otimes x^{(n)}$.

First assume that $\tilde{q} \neq \tilde{t}$. Then there is an index $i$
such that $x^{(i)} \neq y^{(i)}$. Since $x^{(i)},y^{(i)}$ are both in the
probability simplex $\Delta_{d^{(i)}-1}$, this implies that
$x^{(i)},y^{(i)}$ are linearly independent. Let
$A$ be the $d^{(i)}_1 \times \left(\prod_{i' \neq i} d_{i'}\right)$-matrix whose
rows are the slices of $p$ labeled by $k \in [d^{(i)}_1]$ (vectorized in
some fixed but arbitrary manner); so $A$ is part of a
flattening of $p$. Let $B$ be the corresponding part of a flattening
of $q$. Then $A$ and $B$ are rank-one matrices with column spaces spanned
by $x^{(i)}$ and $y^{(i)}$, respectively, hence different column
spaces since $x^{(i)},y^{(i)}$ are linearly independent.

Now at least one of $i_1,i_2$ from item (3) above is not equal to $i$;
say that $i_2 \neq i$. This implies that the rows of $A$ have positive
entries in positions whose $i_2$-th coordinate is $>d^{(i_2)}_1$,
whereas $B$ only has zero entries there. This means that also the column
spaces of $A$ and $B$ are distinct.  Then for any $0 < \varepsilon <1$
the combination $C:=(1-\varepsilon)A + \varepsilon B$ is a rank-two matrix. 
Moreover, by Lemma \ref{LemmaSubgameAdd}, we may choose this
$\varepsilon$ such that $p':=(1-\varepsilon)p+\varepsilon q$ lies in $P_G$.
We conclude that the slices of $p'$ in the $i$-th direction labeled by
$k \in [d^{(i)}_1]$, which correspond to the rows of $C$, are not all
scalar multiples of each other. Now, Remark \ref{RemarkR}(3) implies that (by potentially decreasing $\varepsilon$ further) $r(p')>r(p)$. 

We are left with the case that $\tilde{q} = \tilde{t}$. Recall that $d^{(1)}_1 \geq 2$; we will adjust the payoff tensor $X^{(1)}$. 
At least one of $i_1$ and $i_2$ is not equal to $1$, say $i_2 \neq 1$. Now 
$p_{-1,1}$ has positive entries on positions whose $i_2$-coordinate
is $>d^{(i_2)}_1$, while $q_{-1,1}$ has zeros there. Hence the 
slices $p_{-1,1}$ and $q_{-1,1}$ are linearly independent, and there
exists a linear function $h \in (\RR^{D/d_i})^*$ with 
$h(p_{-1,1})=0>h(q_{-1,1})$. Define the
$(d_1 \times \dots \times d_n)$-tensor $Y^{(1)}$ by setting its $k$-th slice
in the first direction to be 
\[Y^{(1)}_k= \begin{cases}
    X^{(1)}_1 + \varepsilon h & \text{if } k=1, \text{ and}\\
    X^{(1)}_k & \text{otherwise.}
\end{cases} 
\] for some small $\varepsilon>0$. Now set $X':=(Y^{(1)},X^{(2)}, \dots,
X^{(n)})$ and $G':=(n,S,X')$ with the subgame $\tilde{G'}$ corresponding
to the subsets strategies $M^{(1)}_1, \dots, M^{(n)}_1$. We claim that, if
$\varepsilon$ is sufficiently small, then $p \in P_{G'}$ and $p$ is balanced
for $G'$. Indeed, the only incentive constraints that may have changed
compared to those for $G$ are the inequalities
$H^{Y^{(1)}}_{k,\ell}(p_{-1,k})
\geq 0$ when at least one of $k,\ell$ is equal to $1$. However:
\begin{itemize} 
\item if $k=1$, then since $h(p_{-1,1})=0$,
the constraint in fact does not change; 
\item if $\ell=1$ and $k \in [d^{(1)}_1] \setminus \{1\}$, then $p_{-1,k}$ is a scalar multiple of
$p_{-1,1}$ and hence $h(p_{-1,k})=0$, and again the constraint does
not change; and 
\item if $\ell=1$ and $k>d^{(1)}_k$, then since $p$ is balanced we have
a strict inequality $H^{X^{(1)}}_{k,\ell}(p_{-1,k})>0$, hence
choosing $\varepsilon$
sufficiently small we still have $H^{Y^{(1)}}_{k,\ell}(p_{-1,k})>0$. 
\end{itemize}
The balancedness of $p$ for $G'$ follows because we keep all the
strict inequalities among the incentive constraints. 

Now the first two slices of $q$ in the first direction are positive
scalar multiples of each other; this follows from 
$d^{(1)}_1>1$, the fact that $q$
has rank $1$, and the fact that $y^{(1)}=x^{(1)}$, where the latter vector
only has positive entries. We therefore have
$H^{X^{(1)}}_{1,2}(q_{-1,1})=0$ by Remark~\ref{re:IncConsEq}. Hence
\[ H^{Y^{(1)}}_{1,2}(q_{-1,1})=H^{X^{(1)}}_{1,2}(q_{-1,1})+\varepsilon h(q_{-1,1})= \varepsilon h(q_{-1,1})<0.\] Thus $\tilde{q} \not\in P_{\tilde{G'}}$, i.e., $(x^{(1)},\dots, x^{(n)})$ is not a Nash equilibrium of $\tilde{G'}$. Finally, let $(z^{(1)}, \dots, z^{(n)})$ be a Nash equilibrium of $\tilde{G'}$ and let $\tilde{u}:=z^{(1)} \otimes \dots \otimes z^{(n)}$. Then $\tilde{u} \neq \tilde{t}$, so we have reduced to the former case and we are done.
\end{proof}

\end{document}
